% The four parts of the Day 9 lab. The steps follow Labs/day09/README.md.
%
% A value that only a board can give is written as "..." and the students
% measure it. The deck does not show a result that the students must predict.

% ------------------------------------------------------------------- Part A
\begin{frame}{Part A: the protocol (20 min)}
  \small
  \renewcommand{\arraystretch}{1.15}
  \begingroup\centering
  \begin{tabular}{@{}L{0.07\textwidth}L{0.70\textwidth}L{0.12\textwidth}@{}}
    \toprule
    \textbf{Step} & \textbf{Work} & \textbf{Time} \\
    \midrule
    1 & Read the plan: three groups of measurements in the tables of Parts B
        and C of \code{report.md} & 5 min \\
    2 & Fill the protocol table of Part A: the seven parts of a benchmark
        for the two boards & 10 min \\
    3 & Write a prediction with a reason for each question of the
        prediction table & 5 min \\
    \bottomrule
  \end{tabular}\par\endgroup
  \vskip 0.4em
  \begin{itemize}
    \item Decide the warm-up runs and the timed runs. The sketch has 3 and
          20 at the start, because one inference on the kit can need more
          than 1 second. The script of the Raspberry Pi has 20 and 200.
    \item The statistics: the median and the 95th percentile. The window:
          the model only. The input: one fixed input for each model.
    \item Do not change the protocol after a result. If you must change
          the method, write the change and the reason.
  \end{itemize}
  \begin{exerciseblock}
    \textbf{Predictions.} Is \code{int8} faster than \code{float32} on the
    kit? How many times faster is the Raspberry Pi for the same file? Which
    board uses less energy for one keyword inference? Give a number and a
    reason for each question.
  \end{exerciseblock}
\end{frame}

% ------------------------------------------------------------------- Part B
\begin{frame}{Part B: the microcontroller (45 min)}
  \small
  \renewcommand{\arraystretch}{1.15}
  \begingroup\centering
  \begin{tabular}{@{}L{0.07\textwidth}L{0.70\textwidth}L{0.12\textwidth}@{}}
    \toprule
    \textbf{Step} & \textbf{Work} & \textbf{Time} \\
    \midrule
    1 & Build and upload \code{sketches/kit\_bench/kit\_bench.ino} & 10 min \\
    2 & Task B1: the function \code{percentileUs} in
        \code{bench\_stats.h} & 10 min \\
    3 & Measure the four models, copy the lines \code{CSV}, measure again
      & 10 min \\
    4 & The clock: \code{CPU\_MHZ} 80, then 240 again & 10 min \\
    5 & Compare with the published results and with Day 5 & 5 min \\
    \bottomrule
  \end{tabular}\par\endgroup
  \vskip 0.4em
  \begin{itemize}
    \item Board \code{XIAO\_ESP32S3}, and \code{Tools > PSRAM > OPI PSRAM}.
          Set \code{WARMUP\_RUNS} and \code{TIMED\_RUNS} at the top of the
          sketch to the values of your protocol.
    \item The first build needs some minutes. Do Task B1 during the build:
          the nearest-rank rule, with integer arithmetic.
    \item \textbf{You record:} the two lines \code{Sketch uses ... bytes}
          and \code{Global variables use ... bytes} of the build.
  \end{itemize}
\end{frame}

\begin{frame}[fragile]{Part B: the output of the sketch}
  \small
  The Serial Monitor at 115\,200 baud shows \code{Task B1: complete}, and
  then for each of the four models:
\begin{lstlisting}
--- kws int8 ---
model: 53936 bytes
arena: ... bytes used of 32768, in the internal RAM
first inference: ... us
timed runs: 20, after 3 inferences to warm up
min ... us, median ... us, p95 ... us, max ... us
output: class 8 (LiteRT: 8), largest difference ... Same result as LiteRT: yes
chip temperature: ... C
CSV,kit,kws,int8,53936,...
\end{lstlisting}
  \begin{itemize}
    \item Make the file \code{results\_kit.csv} in \code{Labs/day09/}. Copy
          each line that starts with \code{CSV} into it.
    \item \code{Same result as LiteRT: yes} checks the model: the kit gives
          the output of the laptop for the same input.
    \item Send a character in the Serial Monitor: the sketch measures again.
          The difference of the two medians is your noise.
    \item \textbf{You record:} for each model the flash, the arena and its
          memory, the first inference, the median, the p95, and the
          maximum.
  \end{itemize}
\end{frame}

\begin{frame}{Part B: the clock and the published results}
  \small
  \textbf{Step 4.} Set \code{CPU\_MHZ} to 80. Upload, and copy the new
  lines \code{CSV} into \code{results\_kit.csv}, after the first lines.
  Then set \code{CPU\_MHZ} to 240 again.
  \vskip 0.4em
  \textbf{Step 5.} The same two \code{int8} reference models in MLPerf Tiny
  v1.2:
  \vskip 0.2em
  \footnotesize
  \renewcommand{\arraystretch}{1.15}
  \begingroup\centering
  \begin{tabular}{@{}L{0.26\textwidth}L{0.36\textwidth}rr@{}}
    \toprule
    \textbf{Board} & \textbf{Processor} & \textbf{Keywords}
    & \textbf{Images} \\
    \midrule
    NUCLEO-L4R5ZI   & Cortex-M4 at 120 MHz  & 43.2 ms & 165.1 ms \\
    NUCLEO-U575ZI-Q & Cortex-M33 at 160 MHz & 29.1 ms & 111.1 ms \\
    NUCLEO-H7A3ZI-Q & Cortex-M7 at 280 MHz  & 11.3 ms & 40.8 ms \\
    Your XIAO ESP32S3 & Xtensa LX7 at 240 MHz & ... & ... \\
    \bottomrule
  \end{tabular}\par\endgroup
  \vskip 0.4em
  \small
  \begin{itemize}
    \item The three boards use an inference tool of their vendor. Your kit
          uses portable kernels of TensorFlow Lite Micro.
    \item If you have your report of Day 5, write the two times of your own
          keyword model. Which parts of today's protocol did Day 5 not
          have?
  \end{itemize}
  \source{MLPerf Tiny v1.2, closed division (MLCommons).\\
    The latency is 1000 divided by the published rate.}
\end{frame}

% ------------------------------------------------------------------- Part C
\begin{frame}[fragile]{Part C: the Raspberry Pi (50 min)}
  \small
  \renewcommand{\arraystretch}{1.02}
  \begingroup\centering
  \begin{tabular}{@{}L{0.07\textwidth}L{0.70\textwidth}L{0.12\textwidth}@{}}
    \toprule
    \textbf{Step} & \textbf{Work} & \textbf{Time} \\
    \midrule
    1 & Copy the folder, connect, and check the files of Days 7 and 8
      & 5 min \\
    2 & Task C1: the function \code{percentile} in \code{pi/bench.py}
      & 5 min \\
    3 & Idle: the temperature and the power with no load & 5 min \\
    4 & The kit models and MobileNetV2 & 10 min \\
    5 & The detector, in the second environment & 10 min \\
    6 & A sustained run of 180 seconds & 10 min \\
    7 & Copy the results to the laptop & 5 min \\
    \bottomrule
  \end{tabular}\par\endgroup
\begin{lstlisting}[style=shell]
scp -r day09 edge@pi-NN.local:~/edgeai/      # on the laptop, in Labs/
ssh edge@pi-NN.local                         # then on the board
cd ~/edgeai/day09
source ~/tflite_env/bin/activate
ls ../day07/models ../day08/models
python pi/bench.py --idle 30
\end{lstlisting}
  \begin{itemize}
    \item You must see \code{mnv2.tflite} and \code{mnv2\_static.onnx} of
          Day 7, and the three \code{cupbottle\_320} files of Day 8.
    \item \textbf{Task C1:} the same rule as Task B1, in Python.
  \end{itemize}
\end{frame}

\begin{frame}[fragile]{Part C: three suites}
\begin{lstlisting}[style=shell]
python pi/bench.py --suite tiny              # the four models of the kit
python pi/bench.py --suite classification    # MobileNetV2 of Day 7
deactivate
source ~/yolo/bin/activate                   # this environment has NCNN
python pi/bench.py --suite detection         # your detector of Day 8
\end{lstlisting}
  \small
  The script prints one line for each case, and adds one row to the file
  \code{results\_pi.csv}:
  \vskip 0.2em
  \footnotesize
  \renewcommand{\arraystretch}{1.1}
  \begingroup\centering
  \begin{tabular}{@{}L{0.30\textwidth}L{0.62\textwidth}@{}}
    \toprule
    \textbf{Column} & \textbf{Content} \\
    \midrule
    task, model, runtime, type & The case: one file with one runtime \\
    threads                    & 1 and 4 (the suite \code{tiny}: 1) \\
    first, median, p95, max    & Times in ms: the first inference alone,
                                 then the timed runs \\
    RSS MB, temp, watt         & Peak RAM of the process, highest
                                 temperature, estimated power \\
    check                      & \code{tiny}: \code{same as the kit sketch} \\
    \bottomrule
  \end{tabular}\par\endgroup
  \vskip 0.3em
  \small
  \begin{itemize}
    \item Each case runs in a new process. Use \code{--warmup} and
          \code{--runs} for other values than 20 and 200.
  \end{itemize}
\end{frame}

\begin{frame}[fragile]{Part C: the sustained run}
\begin{lstlisting}[style=shell]
python pi/bench.py --sustain 180 --threads 4 \
    --model ../day08/models/cupbottle_320.tflite
\end{lstlisting}
\begin{lstlisting}
  second inferences  median ms     p95 ms   temp C clock MHz  throttled    watt
      10        ...        ...        ...      ...       ...        ...     ...
      ...
     180        ...        ...        ...      ...       ...        ...     ...
Median latency: ... ms in the first window, ... ms in the last window (...).
\end{lstlisting}
  \small
  \begin{itemize}
    \item One line for each 10 seconds. The script writes the lines into
          \code{sustain.csv}.
    \item \code{throttled}: \code{0x0} means no event since the start of the
          board. Part 3 of the lecture gives the other values.
    \item \textbf{You record:} the first window and the last window: the
          median, the p95, the temperature, the clock, and the power.
    \item Then copy \code{results\_pi.csv} and \code{sustain.csv} to the
          laptop with \code{scp} (step 7 of the README).
  \end{itemize}
  \begin{exerciseblock}
    \textbf{Prediction.} Does the latency of the last window differ from
    the latency of the first window? Write your answer before the run.
  \end{exerciseblock}
\end{frame}

% ------------------------------------------------------------------- Part D
\begin{frame}{Part D: the report (35 min)}
  \small
  \renewcommand{\arraystretch}{1.15}
  \begingroup\centering
  \begin{tabular}{@{}L{0.07\textwidth}L{0.70\textwidth}L{0.12\textwidth}@{}}
    \toprule
    \textbf{Step} & \textbf{Work} & \textbf{Time} \\
    \midrule
    1 & The power of each board in \code{power.csv} & 10 min \\
    2 & Tasks D1 and D2 in \code{pi/make\_report.py} & 5 min \\
    3 & Run \code{python3 pi/make\_report.py} on the laptop & 5 min \\
    4 & The accuracy table, the answers, and the Decision Log & 15 min \\
    \bottomrule
  \end{tabular}\par\endgroup
  \vskip 0.4em
  \begin{itemize}
    \item \textbf{Power meter:} read the board with no load and during a
          run. Write the value and the method in \code{power.csv}.
    \item \textbf{No meter:} keep the data sheet values for the kit (217,
          157, and 0.79 mW). For the Raspberry Pi, keep the row
          \code{pi,active} empty: the script uses the estimate of the board.
    \item \textbf{Task D1:} energy = power $\times$ latency. \textbf{Task
          D2:} the mean power of a duty cycle.
    \item The script prints four tables and writes them into
          \code{report\_tables.md}: the kit, the Raspberry Pi, the same file
          on the two boards, and a battery estimate.
  \end{itemize}
  \keypoint{An energy value needs its boundary and its instrument: the chip
    or the board, the data sheet or the meter.}
\end{frame}

\begin{frame}{Part D: the accuracy of each model}
  \small
  You do not measure the accuracy today. Use these values, with their
  source:
  \vskip 0.3em
  \footnotesize
  \renewcommand{\arraystretch}{1.15}
  \begingroup\centering
  \begin{tabular}{@{}L{0.33\textwidth}L{0.15\textwidth}L{0.42\textwidth}@{}}
    \toprule
    \textbf{Model} & \textbf{Accuracy} & \textbf{Test set and source} \\
    \midrule
    Image classification, \code{float32} & 87.2 percent
        & 10 000 test images of CIFAR-10, this course \\
    Image classification, \code{int8}    & 87.0 percent
        & 10 000 test images of CIFAR-10, this course \\
    Keyword spotting, \code{int8}        & 91.6 percent
        & Speech Commands, published by MLCommons \\
    Keyword spotting, \code{float32}     & not measured
        & \code{models/README.md} \\
    MobileNetV2: \code{float32}, \code{int8} & 77.4 and 75.2 percent
        & 500 images of Imagenette, the Day 7 lab \\
    Your detector                        & your values
        & Your notebook of Day 8 \\
    \bottomrule
  \end{tabular}\par\endgroup
  \vskip 0.4em
  \small
  \begin{itemize}
    \item The two image files give the same accuracy on the 200 test images
          of MLPerf Tiny: 87.0 percent. The interval of 200 images is plus
          and minus 4.7 points.
    \item A latency with no accuracy can reward the wrong file. Each row of
          your report needs the two numbers.
  \end{itemize}
  \source{Image classification: an experiment of this course.\\
    Keyword spotting: Banbury et al., 2021.}
\end{frame}
